\section{A global higher-product family with embedded structure}
\label{sec:higher-hyperplane-classification}
We now apply the residual mechanism to a natural family in every
symmetric degree, determining the whole reduced failure locus rather
than only an obstruction subfamily. The scheme structure contains a
further phenomenon: its support is independent of the degree, but
its nilpotent structure is not.

Fix $n\ge4$ and put $V_n=H^0(\Pj^1,\mathcal O(n))$. A hyperplane
$U\subset V_n$ is the kernel of a unique point $[h]\in\Pj(V_n^*)$.
For $m\ge2$ let
\[
 \mu_m:\Sym^m\mathcal U\longrightarrow
          V_{mn}\otimes\mathcal O_{\Pj(V_n^*)},\qquad
 D_m=V(I_{mn+1}(\mu_m)).
\]
Its source rank is $e_m=\binom{n+m-1}{m}\ge mn+1$.
Write $C_n\subset\Pj(V_n^*)$ for the evaluation rational normal
curve, and $S_n$ for its secant threefold, including tangent secants.
Let $\mathcal J_m$ be the ideal sheaf of $D_m$, and
$\mathcal I_{S_n},\mathcal I_{C_n}$ the reduced ideal sheaves.
Associated points below are those of coherent sheaves on projective
space; no assertion about an irrelevant-ideal component of an affine
homogeneous presentation is implicit.

\begin{theorem}[Higher-product hyperplanes: support, ranks and associated points]
\label{thm:higher-hyperplane-global}
For every $n\ge4$ and $m\ge2$,
\begin{equation}\label{eq:hyperplane-complete-coranks}
 (D_m)_{\rm red}=S_n,\qquad
 \operatorname{corank}\mu_{m,\ker h}=
 \begin{cases}
 0,&[h]\notin S_n,\\
 1,&[h]\in S_n\setminus C_n,\\
 m,&[h]\in C_n.
 \end{cases}
\end{equation}
The scheme $D_m$ is smooth of dimension three at every point of
$S_n\setminus C_n$, including the non-evaluation tangent secants.
In particular
\begin{equation}\label{eq:hyperplane-saturation}
 \mathcal J_m:\mathcal I_{C_n}^{\infty}=\mathcal I_{S_n}.
\end{equation}
For every $m\ge3$, the complete set of associated points is
\begin{equation}\label{eq:hyperplane-associated-points}
 \operatorname{Ass}_{\mathcal O_{\Pj(V_n^*)}}
                  (\mathcal O_{D_m})=\{\eta_{S_n},\eta_{C_n}\}.
\end{equation}
Thus there is exactly one reduced irreducible component, and exactly
one embedded associated subvariety, the evaluation curve. Its
nilradical
$\mathcal N_m=\mathcal I_{S_n}/\mathcal J_m$ satisfies, at every
closed point $x\in C_n$,
\begin{equation}\label{eq:hyperplane-nilpotence-bound}
 (\mathcal N_m^j)_x\ne0\quad\text{whenever }1\le j\text{ and }2j<m.
\end{equation}
Consequently the nilpotency index at $x$ is at least
$\lceil m/2\rceil$. The theorem does not assert that this lower
bound is the exact index.
\end{theorem}

\begin{lemma}[Propagation from quadratic surjectivity]
\label{lem:hyperplane-propagation}
Let $U\subset V_n$ be base-point-free. Then
$UV_d=V_{n+d}$ for every $d\ge n-1$. In particular,
$U^2=V_{2n}$ implies $U^m=V_{mn}$ for every $m\ge2$.
\end{lemma}
\begin{proof}
There are coprime $f,g\in U$: fix a nonzero $f$ and choose $g$
outside the finitely many proper evaluation hyperplanes at its
zeros. The kernel of
\[
 V_d\oplus V_d\longrightarrow V_{n+d},\qquad
                  (A,B)\longmapsto fA+gB
\]
consists exactly of $(gh,-fh)$ with $h\in V_{d-n}$.
For $d=n-1$ its kernel is zero and both dimensions are $2n$.
For $d\ge n$ the rank is
$2(d+1)-(d-n+1)=d+n+1$. Thus the map is onto in both cases.
If $U^2=V_{2n}$ then $U$ is automatically base-point-free;
apply the first assertion successively with $d=2n,3n,\ldots$.
\end{proof}

\begin{lemma}[All higher-product ranks on secants]
\label{lem:hyperplane-secant-ranks}
For a non-evaluation point $h\in S_n$, the image $U^m$ has
codimension one in $V_{mn}$ for every $m\ge2$. For an evaluation
point it has codimension $m$.
\end{lemma}
\begin{proof}
First take a distinct-support secant and choose the two points as
$0$ and $\infty$. In the binary monomial basis
$z_j=X^jY^{n-j}$, its hyperplane has a basis
\[
 z_1,\ldots,z_{n-1},\quad a=z_0+\lambda z_n,
                       \qquad\lambda\ne0.
\]
Products of two of the interior monomials give all exponents
$2,\ldots,2n-2$. Multiplication of $a$ by $z_1$ and $z_{n-1}$
then supplies exponents $1$ and $2n-1$, since the other terms are
already in that span. Finally $a^2$ gives the one allowed combination
of the endpoint monomials. Thus $U^2$ contains every interior
monomial and one endpoint combination, and has dimension $2n$.

Inductively suppose $U^m$ contains all monomials of exponents
$1,\ldots,mn-1$, as well as $a^m$. Multiplying its interior
monomials by $z_1,\ldots,z_{n-1}$ gives exponents
$2,\ldots,(m+1)n-2$. Multiplying its two extreme interior monomials
by $a$ supplies the two missing interior exponents. The element
$a^{m+1}$ supplies an endpoint combination. All products satisfy the
single endpoint relation inherited from $f(0)=\lambda^{-1}f(\infty)$
for $f\in U$. Hence there can be no second endpoint direction, and
the image has dimension $(m+1)n$.

At a non-evaluation tangent secant, a change of homogeneous
coordinates puts the hyperplane in the form with the first derivative
at $0$ equal to zero. To justify the normalization, in a local
coordinate a change fixing $0$ changes a first-jet functional by a
nonzero scalar and an arbitrary multiple of evaluation; the latter
multiple can therefore be removed. The hyperplane is then
\[
 U=\operatorname{span}\{z_0,z_2,z_3,\ldots,z_n\}.
\]
Sums of $m$ exponents in $\{0,2,3,\ldots,n\}$ give exactly
$\{0,2,3,\ldots,mn\}$. This follows by induction, since the
intervals obtained by adding $[2,n]$ overlap or are adjacent for
$n\ge3$. Thus precisely exponent $1$ is missing.

Finally, for $h$ an evaluation, $U=lV_{n-1}$ for its vanishing
linear form $l$. Since products of the complete binary series span,
$U^m=l^mV_{m(n-1)}$, of dimension $m(n-1)+1$. Its corank is $m$.
\end{proof}

\begin{lemma}[Explicit plane-derivative matrix and residual span]
\label{lem:hyperplane-explicit-matrix}
Use the monomials $z_j=X^jY^{n-j}$. At a split secant take
\[
 a=z_0+\lambda z_n,\quad (w_1,\ldots,w_{n-1})=(z_1,\ldots,z_{n-1}),
 \quad v=z_n,\quad\lambda\ne0,
\]
and let the annihilator be coefficient $mn$ minus $\lambda^m$
times coefficient zero. At a tangent secant take $a=z_0$,
$(w_1,\ldots,w_{n-1})=(z_2,\ldots,z_n)$, $v=z_1$, and the
coefficient-one annihilator. Set $u_0=a$, $u_i=w_i$, and vary
$u_i$ by $t_i v$. On the ordered source quotient
\[
 a^m,w_1a^{m-1},\ldots,w_{n-1}a^{m-1}
\]
the plane-derivative matrix is
\begin{equation}\label{eq:hyperplane-explicit-diagonal}
       \kappa\,\operatorname{diag}(m,1,\ldots,1),\qquad
       \kappa=\lambda^{m-1}\ \text{or}\ 1,
\end{equation}
respectively. It is zero on all monomials with at least two $w$
factors. Their product image has dimension $mn-3$ in both cases.
\end{lemma}
\begin{proof}
Let $\ell$ be the indicated annihilator. Then
$\ell(va^{m-1})=\kappa\ne0$. The variation of $a^m$ gives
$m\kappa t_0$. For $w_i a^{m-1}$ only variation of $w_i$ can
contribute to $\ell$, giving $\kappa t_i$; all terms which retain
$w_i$ vanish. For a monomial with at least two $w$ factors one
such factor remains after every single variation, and the result
is annihilated by $\ell$. This gives every entry of
\eqref{eq:hyperplane-explicit-diagonal}.

In the split case $W=XYV_{n-2}$, so
$W^2=(XY)^2V_{2n-4}$. The subseries $U=W+\C a$ is
base-point-free: $W$ has only the two endpoint base points, and
$a$ is nonzero at both. Since $2n-4\ge n-1$ for $n\ge4$,
Lemma~\ref{lem:hyperplane-propagation} can be applied successively
at degrees $2n-4,3n-4,\ldots$, proving
$W^2U^{m-2}=(XY)^2V_{mn-4}$.

In the tangent case the exponent set of $W^2$ is the entire
integer interval $[4,2n]$. If the exponent set at degree $k$ is
$[4,kn]$, adding the exponent set $\{0\}\cup[2,n]$ of $U$
gives $[4,kn]\cup[6,(k+1)n]=[4,(k+1)n]$, since $kn\ge5$.
Induction from $k=2$ proves the assertion in every degree.
There are exactly $mn-3$ exponents in $[4,mn]$.
\end{proof}


\begin{lemma}[Scheme smoothness off the evaluation curve]
\label{lem:hyperplane-residual-smoothness}
At every point of $S_n\setminus C_n$, the plane derivative of the
annihilator incidence has rank $n$, and its polar residual image has
dimension $mn-3$. Consequently the tangent space of $D_m$ has
dimension three at each such point.
\end{lemma}
\begin{proof}
For distinct support, let $g$ be its vanishing quadratic and write
$U=W\oplus\C a$, with $W=gV_{n-2}$. The plane derivative vanishes
on monomials with at least two factors in $W$, while it is an
isomorphism onto the dual of the $n$-dimensional quotient spanned by
$a^m$ and $wa^{m-1}$. Indeed the functional on $V_n/U$ obtained by
varying one factor and evaluating against the annihilator is nonzero;
the images of the $n$ basis vectors of $U$ can be varied independently.
This is the computation of Theorem~\ref{thm:higher-residual}, now
with $c=n$. The residual image is
\[
 W^2U^{m-2}=g^2V_{2n-4}U^{m-2}=g^2V_{mn-4}.
\]
The last equality follows from
Lemma~\ref{lem:hyperplane-propagation}, because $U$ is
base-point-free and $2n-4\ge n-1$. Its dimension is $mn-3$.

For a tangent secant in the normal form of
Lemma~\ref{lem:hyperplane-secant-ranks}, let
$W=\operatorname{span}\{z_2,\ldots,z_n\}$ and $a=z_0$.
The annihilator in $V_{mn}^*$ is the coefficient of exponent $1$.
Since $V_n/U$ is spanned by $z_1$, the derivative independently
prescribes the coefficients dual to $a^m$ and $wa^{m-1}$, and has
rank $n$. Once again its kernel on the dual source consists of
monomials with at least two $W$ factors. Their products span precisely
the monomials of exponents $4,\ldots,mn$: $W^2$ gives the full
interval $[4,2n]$, and successive factors from $U$ extend it without
gaps. This image also has dimension $mn-3$.

The incidence tangent sequence now gives, in both cases,
\[
 \dim T D_m=\dim T\mathcal I
      =n+(mn+1)-1-n-(mn-3)=3,
\]
because the original multiplication has corank one by
Lemma~\ref{lem:hyperplane-secant-ranks}.
\end{proof}

\begin{lemma}[The length-two radical, including a double point]
\label{lem:hyperplane-confluent-radical}
Let $0\ne\ell\in V_{2n}^*$ and let
$H_\ell(f,g)=\ell(fg)$ on $V_n$. If $H_\ell$ has rank two,
there is a length-two divisor $Z\subset\Pj^1$ such that
\[
 \ell\in H^0(\mathcal O_Z(2n))^*\subset V_{2n}^*,\qquad
 \operatorname{rad}H_\ell=H^0(\mathcal O(n)(-Z)).
\]
The induced pairing on $H^0(\mathcal O_Z(n))$ is nondegenerate.
\end{lemma}
\begin{proof}
The rank-at-most-two Hankel variety is the secant variety of the
evaluation rational normal curve. The shape-independence of the
Hankel minor ideal and the precise secant parametrization are given
in \cite[Section 1 and Section 2, equation (2.0.1)]{CMSV}; here the
minor size is three. The projective bundle of the spaces
$H^0(\mathcal O_Z(2n))^*$ over $\operatorname{Hilb}^2\Pj^1$ is
proper. Its image contains the distinct-point secant lines and is
contained in the rank-at-most-two locus by factorization through a
two-dimensional restriction space. Its image is therefore the
entire secant variety, including tangent lines. Thus $\ell$
factors through some $Z$.

For two distinct points the pairing on the length-two algebra has
matrix $\operatorname{diag}(\alpha,\beta)$. Rank two forces
$\alpha\beta\ne0$. For $Z=2p$, choose a parameter $z$ and
trivialize the line bundle. Write
$\ell(A+Bz)=\alpha A+\beta B$ on $\C[z]/(z^2)$. The pairing
matrix in the basis $1,z$ is
$\left(\begin{smallmatrix}\alpha&\beta\\\beta&0\end{smallmatrix}\right)$,
which has rank two exactly when $\beta\ne0$. Restriction
$V_n\to H^0(\mathcal O_Z(n))$ is surjective. In either case the
kernel of the pulled-back nondegenerate pairing is exactly its
restriction kernel $H^0(\mathcal O(n)(-Z))$.
\end{proof}

\begin{lemma}[Equivariant associated points on a transitive curve]
\label{lem:equivariant-associated-curve}
Let a group act on a Noetherian variety, preserving an ideal
$\mathcal J$ with irreducible reduced support $S$. Let $C\subset S$
be a closed irreducible curve on whose closed points the group acts
transitively. Suppose $\mathcal J=\mathcal I_S$ away from $C$
and the nonzero module $\mathcal N=\mathcal I_S/\mathcal J$
has support exactly $C$. Then
\[
       \operatorname{Ass}(\mathcal O/\mathcal J)
                     =\{\eta_S,\eta_C\}.
\]
\end{lemma}
\begin{proof}
For a finite module over a Noetherian ring, associated primes are
finite, minimal points of its support are associated, and the
associated primes of a submodule are associated to the ambient
module; see \cite[Lemma 10.63.3, Lemma 10.63.5,
and Proposition 10.63.6]{StacksAssociated}. These assertions apply
on a finite affine cover and give the corresponding coherent-sheaf
statements. Associated points outside $C$ consist only of $\eta_S$,
since the quotient there is the structure sheaf of a reduced
irreducible variety. An associated point with closure inside $C$
is either $\eta_C$ or a closed point. The finite associated set is
preserved by the group action. Transitivity on the infinitely many
closed points of $C$ excludes the latter possibility. Finally
$\eta_C$ is a minimal point of the support of $\mathcal N$, hence
is associated to $\mathcal N$ and therefore to
$\mathcal O/\mathcal J$. The minimal point $\eta_S$ is associated
as well. This gives exactly the two asserted points.
\end{proof}

The coordinate assertion in
Lemma~\ref{lem:hyperplane-residual-smoothness} is supplied by
\eqref{eq:hyperplane-explicit-diagonal}. The length-two assertion
and the associated-point step of the following proof are, respectively,
Lemma~\ref{lem:hyperplane-confluent-radical} and
Lemma~\ref{lem:equivariant-associated-curve}. We retain the full
argument to specify where its nilpotence estimate is used.


\begin{proof}[Proof of Theorem~\ref{thm:higher-hyperplane-global}]
We first identify the quadratic support. A hyperplane $U$ can be
totally isotropic for a nonzero Hankel form $H_\ell$ only if that
form has rank at most two: the maximal isotropic dimension of a
rank-$r$ form on an $(n+1)$-space is
$n+1-\lceil r/2\rceil$. A rank-one Hankel form gives an evaluation
hyperplane. A rank-two Hankel form has radical
$H^0(\mathcal O(n)(-Z))$ for a length-two divisor $Z$ on $\Pj^1$.
For clarity, this last assertion includes a double point: the proper
family of secant lines indexed by length-two divisors covers the
rank-at-most-two Hankel variety of Lemma~\ref{lem:hankel-strata};
on a length-two quotient the rank-two pairing is nondegenerate,
so its kernel is exactly the indicated space. Since an isotropic
hyperplane contains this radical, its annihilating functional lies
on the corresponding secant or tangent line to $C_n$. Hence
$|D_2|\subset S_n$. Conversely a distinct-support secant imposes a
nonzero quadratic endpoint relation. Closedness supplies the whole
$S_n$, proving $|D_2|=S_n$.

Outside $S_n$, quadratic multiplication is onto; the propagation
lemma proves surjectivity in every higher degree. On $S_n$ the rank
lemma gives the two remaining cases of
\eqref{eq:hyperplane-complete-coranks}. Thus $|D_m|=S_n$ globally.
Since $S_n$ is an irreducible threefold,
Lemma~\ref{lem:hyperplane-residual-smoothness} equates local Krull
and tangent dimensions at every point off $C_n$. The local rings
there are regular. In particular
$\mathcal J_m=\mathcal I_{S_n}$ away from $C_n$.

The coherent quotient
$\mathcal N_m=\mathcal I_{S_n}/\mathcal J_m$ is therefore
supported on $C_n$, and some power of $\mathcal I_{C_n}$ kills it.
This gives the inclusion of $\mathcal I_{S_n}$ in the left side
of \eqref{eq:hyperplane-saturation}. The reverse inclusion holds
because $\mathcal I_{S_n}$ is prime and $S_n$ is not contained in
$C_n$. This proves the saturation identity.

At an evaluation point $x$, the original corank is $m$. Schur
elimination consequently expresses $\mathcal J_{m,x}$ as the ideal
of $m$-minors of a matrix whose entries vanish at $x$. In the regular
local ring $\mathcal O_{\Pj(V_n^*),x}$, with maximal ideal
$\mathfrak m_x$, one has
\begin{equation}\label{eq:hyperplane-minor-order}
 \mathcal J_{m,x}\subset\mathfrak m_x^m.
\end{equation}
Choose coordinates $h_0,\ldots,h_n$ dual to the binary monomials
so that $x=[1:0:\cdots:0]$, and work on $h_0=1$. The function
\[
 F=\det\begin{pmatrix}1&h_1&h_2\\h_1&h_2&h_3\\h_2&h_3&h_4\end{pmatrix}
\]
vanishes on $S_n$, since every two-evaluation moment sequence has
Hankel rank at most two. Its initial form at $x$ is
$h_2h_4-h_3^2$, a nonzero quadratic. Therefore $F^j$ has order
$2j$, and \eqref{eq:hyperplane-minor-order} implies that its class
modulo $\mathcal J_m$ is nonzero whenever $2j<m$.
Transitivity of $\operatorname{SL}_2(\C)$ on $C_n$ proves
\eqref{eq:hyperplane-nilpotence-bound} at every evaluation point.
For $m\ge3$ this shows in particular that $\mathcal N_m$ is nonzero
with support all of $C_n$.

It remains to exclude any additional associated points. The
multiplication and its maximal-minor ideal are
$\operatorname{SL}_2(\C)$-equivariant. The associated points of a
coherent sheaf form a finite set, stable under automorphisms.
Off $C_n$ the scheme is the smooth irreducible open part of $S_n$,
so its only associated point there is $\eta_{S_n}$. An associated
point whose closure is contained in $C_n$ is either $\eta_{C_n}$
or a closed point of $C_n$. The latter is impossible: transitivity
would force every closed point of the curve into this finite set.
Thus the only possible embedded associated point is $\eta_{C_n}$.
It actually occurs because the nonzero submodule $\mathcal N_m$
has support $C_n$, whose generic point is a minimal point of its
support and hence an associated point. Associated points of a
submodule are associated points of the ambient module. This proves
\eqref{eq:hyperplane-associated-points}, and completes the proof.
\end{proof}

The expected determinantal codimension here is
$\binom{n+m-1}{m}-mn$, while the actual codimension is $n-3$.
For $m\ge3$ this is an excess family with a complete associated-point
classification. The nilpotency bound is unbounded as $m$ grows even
though the reduced threefold and its unique embedded support do not
change. This conclusion is stronger than the two-point-sector
calculation: it excludes every other reduced failure component and
every other embedded associated subvariety in this entire family.
It does not identify the full embedded primary ideal, or assert an
all-contact, all-$c$ higher-product classification.
